\section{Tour Guide through the Ruins}

\begin{quotex}
The intelligent will converge to a common worldview, but the ignorant believe all manner of things. \flright{\textsc{Old wives' tale}}

Justice means minding one's own business and not meddling with other men's concerns. \flright{\textsc{Plato}}

The end of thinking\footnote{\url{https://gornahoor.net/?p=2369}} is called being empty-headed. Many are like that. Actually, the world would improve if more people stopped thinking altogether. \flright{\textsc{The gypsy lady with a gold-capped tooth}}

\end{quotex}
It seems that, especially during trying times, that ``public" intellectuals feel obliged to proffer profound sounding opinions on current events. The topics include self-alienation, social isolation, and so on: usually plausible rearrangements of ready-made ideas. These get repeated on social media due primarily to the alleged reputations of the writers, rather for their actual content. Moreover, it is always an obscure ``they" who are responsible. That is the opposite of the noble attitude:

\begin{quotex}
Nobility, the quality which I call true aristocracy, requires of a man the recognition of his guilt. \emph{The sense that one is being continually affronted is on the other hand, precisely a plebeian feeling}. \flright{\textsc{Nicolai Berdyaev}} 

\end{quotex}
\paragraph{Why Consciousness}
The idea of alienation has a long past, so it has hardly made its appearance in the modern world. It was the showpiece of Karl Marx, but even the ancient Gnostics understood it well. You can lament about nature red in tooth and claw, natural disasters, or man as wolf to man. The Gnostics experienced that in its depth and concluded that only an evil demiurge could have created this world. What makes it so horrific is that man is fully conscious of such horrors, although his consciousness is entirely passive: an observer, yet impotent to do anything about it. What then is the point?

Of course, those influenced by ``science" believe the same thing, don't they? Their claim is likewise that consciousness is a mere epiphenomenon with no causal efficacy. Hence, did it arise just to be acutely aware of pain and suffering? Unlike the Gnostics, the new intellectuals continue on with the bourgeois existence, without taking their own theories seriously.

\emph{There is quite an abyss between cerebral knowledge and knowing it in your lifeblood}.

\paragraph{Why Disasters}
A popular lament is, ``why does God permit evil in the world?" Then the semi-intelligent will conclude that God is either not all good or not all powerful, usually with a smug grin on their face. Actually, we have known for quite some time how and why death entered into the world.

An interesting thought experience would be to design a world in which free beings exist. I am not considering a so-called ``computer simulation", which would be nothing other than the demiurgic world recognized by the Gnostics; I mean a world in which there are actually free beings. It might not be so different from our world if you do it right. In a totality, everything is necessarily interrelated; that is the clue.

There are two considerations: (1) natural evil caused by physical events and (2) moral evils brought about by humans. Death often comes by surprise. Even the Buddhists know it since they claim, ``The fact of death is certain, the time of death is unknown." Hence there is no reason to conclude that mass deaths in a natural disaster is qualitatively different from a fatal motorcycle accident of an individual. The world operates by natural laws, which are discoverable. The challenge is to build better technology, such as improved building standards in areas prone to earthquakes or hurricanes.

Do you really want God to prevent moral evils? For example, suppose you log into your favourite webcam girl and start fapping. Suddenly, the God squad bursts into the room and zips up your pants. Or you fudge your timesheet at the office and Jiminy Cricket steps in and re-enters the correct numbers. Or you get wired up to receive an electric shock to your brain whenever you fantasize about your neighbour's wife.

I doubt that program would last more than a day before widespread protests erupted. Men are quite content to continue in their sinful ways while simultaneously blaming God for allowing evil.

\paragraph{Gift Giving}
Gift giving is a traditional practice. Not so long ago, it was common to bring a dessert or wine when invited to someone's house. Men today seem to be less conscious of that tradition. Nevertheless, it is good practice to offer a gift when requesting time or service from someone, although it is not mandatory, nor even a condition for that service. Yet, if you ask your buddies to help you move, don't you spring for lunch and have cold beers ready for them afterwards?

Obviously, thanks be to God, I don't need a web presence to pay my bills and hope it remains that way. In any event, it would be impossible given the low volume of readers. And I don't need to appeal for ``tips" like some waitress in a cocktail bar. I can afford to keep the web site online for the foreseeable future, but that is not the point.

To prove me wrong, more than a few people have asked me to restore the gift button. Many of them have been overly generous, for which I am very thankful.

Someone recently told me that most of Gornahoor's pageviews are from bots. In that case, I plan to program my own bots in Python; perhaps I'll get better comments that way. (You don't see the rejects.)

\paragraph{The Almost Life}
In the second volume of My Brilliant Friend, Lila is stuck in a Neapolitan ghetto with a husband she does not love, a job she hates, and no friends. Elena has been able to escape by following all the rules, making the right connections, and not making waves. Her life is safe and predictable.

The contrast with her best friend Lila is jarring. Elena suddenly realizes has been living a ``quasi", an ``almost" life, falling short of full commitment. She is on a path to conventional success. Lila, on the other hands, knows no bounds, acts out her emotions with abandon, and pushes things past socially acceptable limits, even to the point of hurting others and especially herself. Nevertheless, for me the idea of an ``almost" life is the more terrifying alternative.

\paragraph{Male Pronouns}
An apparently well regarded female theologian\footnote{\url{https://womenintheology.org/2016/02/16/theologienne-theologien-theologian-from-angelique-arnauld-to-mary-daly/}} was plagued by this profound observation:

\begin{quotex}
What is the significance of saying that Christ came down from heaven ``for us men and for our salvation"? 

\end{quotex}
She seems to believe that it is anti-woman, when it is actually anti-man. Women get their very own pronouns, but men have to share theirs. The change to ``for us women" deliberately excludes men, but as it stands, it includes females as well as men. That is why the grammatically challenged practice of using an explicitly female pronoun where the male pronoun is grammatically correct. (e.g., ``to each his own") This professor doesn't understand that grammatical gender is different from biological sex. Speakers of Romance languages should probably grasp that idea better.

The consequences are immense. That is why we can speak of ``women's rights" but not of ``men's rights". The former conception definitively indicates females, but the latter is ambiguous: does it mean rights for males alone or for all humans, of whatever sex? This is so obvious that even professors can't grasp it.

\paragraph{Love at First Sight}
I watched Stagecoach (1939) a few nights ago, starring \textbf{John Wayne}. The banker sounded like a guest on Fox News: His opening statement was, ``What's good for banks is good for America", he then complained about the National debt then at \$1 billion, and concluded, ``we should have a businessman as president."

In the movie, the banker got arrested at the end, unlike real life, in which they get bailed out.

John Wayne proposed marriage to a woman he had just met on the stagecoach ride. (i.e. What romance was like before Tinder.) Love at first sight was common in old movies. That actually is rather easy, but love six months later is a much different movie.

\paragraph{Loving your Tradition}
Although we believe in the superiority of our Tradition, that is not intended in a triumphalist ways nor doe we necessarily need to promote or try to demonstrate it. We have readers from around the world with different beliefs. Rather, we mean it in this sense:

If I love my wife, I will not covet your wife. That is your business. Furthermore, if I love my wife, I will not say to her, ``Honey, I love you but there are probably 5 or 10 other women who could have made me equally happy." Not at all, I will tell her, ``Honey, you are the only one for me," and mean it.

Far too many people presume that the Church simply refers to the practices, deeds, and misdeeds of a particular terrestrial organisation. That is far, quite far, from the truth. The Church in her fullness has three branches: the Church Militant (earth), the Church Suffering (purgatory), and the Church Triumphant (Heaven)\footnote{\url{https://diocesan.com/church-suffering-militant-triumphant/}}.

On earth the Church Militant includes the apostles, martyrs, theologians, saints, artists, poets, doctors, philosophers, and so on. Those represent the sources for our worldview. That is why we are little disturbed by the antics of any particular pope, bishop, or priest.

The dogma that the gates of hell shall not prevail against the Church must be understood in her threefold fullness, even though buildings may fall and men might fail.

\paragraph{Choosing the Team}
When I was a young boy, we would organise baseball games on our own, without adult supervision. The two best players, would choose the teams, each one in turn selecting one of the boys. The better players would be chosen first, and the least competent would be last. Competence was the sole criterion; there was no favouritism, no one could bribe his way onto a team, and diversity was unheard of. This was no feeling of injustice in that process, since we all knew our ranks. Looking at this a day later, perhaps there was resentment; they turned into the Boomer socialists.

During the worldwide lockdowns, we've learned that we can live just fine without professional sports, but not the workers at power plants. As captain of a team, as if your life depended on it, which of these would you choose first: a ``news" channels or a meat packing plant?

Given the collapse of many economies around the world, we can see who is necessary and who is not. It is clear we can do without professional athletes, talking heads newscasters, and so on. So, were I a team captain, I would choose the power plant workers, farmers and ranchers, truck drivers, meat packing plant workers, and so on. Athletes, actors, Broadway musicals, and so on would be left on the sidelines.

\paragraph{Family Values}
Talk today about family values is completely misguided, as it has degenerated to nothing more than a so-called nuclear family with bourgeois values. It is all exterior, in appearance only. The character and achievements of the family members are barely taken into consideration, so one wonders exactly what values are meant.

Traditionally, family values were understood as multi-generational. What has the family contributed to the wider society? At the top are the rare families who produced statesmen, bishops, artists, scientists, musicians, and so on. Nevertheless, it is not too late to work on your own family values, however modest. First, know what is really of value.

The ancients used to keep a hearth going for their ancestors. Each family in the city had its own rites and rituals. This is called ``honour thy father and thy mother". I'm amazed at the number of people, even self-described ``traditionalists", who reject the life that their ancestors led for the past thousand or more years. Why do they believe their ancestors were deluded, unintelligent, etc.? Like that woman who called the mother who loved her a ``sexist".

\paragraph{Mystery of the Crucifixion}
It is impossible not to take notice of the ignorance, lack of intelligence, cupidity, laziness, and even maliciousness in public figures. The party system in the USA is unsustainable, since power and winning are the only goals, with no bounds set on means. A disagreement on policy is understandable, but not the wickedness that should be obvious, but is not to partisans. Satan is the father of lies, but too many people prefer the lies.

I suppose I am describing the ruins no better than anyone else. However, you would never hear that on the nightly news.

I know partisans who believe that all the good people in the country become democrats, and the bad ones, are republicans. Certainly that cannot occur by chance, and I don't see any evidence of Maxwell's demon making the choice.

If you could detach enough and take a high-level view, you might also see souls falling into hell like snowflakes in a winter storm. You don't believe that to be possible, even though every traditional religion teaches it? You could take the secular view, but that leaves you with no skin in the game. And that is true for a few professions like journalism, law or politics, in which no one pays a price for being wrong. Everyone else will suffer for being so wrong so often.

But here is the mystery that bothers me. There are a few precious people whom I would die for, but absolutely not for those people just described. So when I hear the Jesus died even for their sins, I am awestruck.



\flrightit{Posted on 2020-05-17 by Cologero }

\begin{center}* * *\end{center}

\begin{footnotesize}\begin{sffamily}



\texttt{Sviatoslav on 2020-05-21 at 06:28 said: }

Very nice article, thank you Cologero. Indeed, we are men among the ruins, though, the light did not completely disappear. Hopefully, this won't be off-topic, I just want to share something positive. I have recently discovered this youtube channel \url{https://www.youtube.com/user/tfpstudentaction/videos} and have been searching a little bit more about that organization and in the end, I have been positively surprised. They are not some neoconservative evangelists as I thought firstly, but truly traditional Catholics. 

At least on videos, I see healthy young men with potential (unlike neopagan nationalists, they don't have tattoos and barbaric look). For example:

\url{https://www.youtube.com/watch?v=gPf6eKfmfZs}

\url{https://tfpstudentaction.org/get-involved/call-to-chivalry-camps/catholic-manhood}


\hfill

\texttt{Sylvan Savant on 2020-05-23 at 19:03 said: }

``Nobility, the quality which I call true aristocracy, requires of a man the recognition of his guilt. The sense that one is being continually affronted is on the other hand, precisely a plebeian feeling."

This statement encapsulates everything of worth and value that has been relayed through this website. Even if it was discussed previously, I still believe it is worth revisiting the profound meaning of these words every now and again. The importance of this message can't be stressed enough.

The world is indeed a hell not because it was created this way, but because WE are responsible for making it thus. The Bible illustrates this fundamental truth on its every page.

Despite being thoroughly of a ``fringe" persuasion, I am becoming increasingly skeptical of the prevailing narrative of the evil elites being responsible for all the world's ills. Don't like how Uncle Gill wants to inject you with his gecko DNA? Stop using Microsoft products. Maybe in a few months or so he will be serving you at a nearby burger joint. Slap his ass for me and don't forget to tip generously.


\end{sffamily}\end{footnotesize}
